% Source problem: 2014ap_exam.pdf, PDF pp. 4--5.
\documentclass[11pt,a4paper]{article}
\usepackage[margin=24mm]{geometry}
\usepackage{fontspec}
\usepackage{xeCJK}
\setCJKmainfont{HaranoAjiMincho-Regular.otf}
\setCJKsansfont{HaranoAjiGothic-Medium.otf}
\usepackage{amsmath,amssymb}
\usepackage{tikz}
\usetikzlibrary{arrows.meta}
\setlength{\parindent}{1em}
\setlength{\parskip}{0.35em}
\title{物理工学専攻 院試解答集\\2014年 物理学I 第2問}
\author{政岡凜太郎}
\date{}
\begin{document}
\maketitle

\section*{[1] 無限直線電荷}
線電荷から距離 $r$ の同軸円筒をガウス面に選ぶ。長さ $l$ の円筒面を貫く電束は $2\pi rl E(r)$ で、内部電荷は $\lambda l$ だから
\begin{equation}
 \boldsymbol E(r)=\frac{\lambda}{2\pi\varepsilon_0r}\,\boldsymbol e_r,
 \qquad E(r)=\frac{|\lambda|}{2\pi\varepsilon_0r}.
\end{equation}
$r_0$ を基準に $E_r=-d\phi/dr$ を積分すると
\begin{equation}
 \phi(r)-\phi(r_0)=-\int_{r_0}^r E_r(r')\,dr'
 =-\frac{\lambda}{2\pi\varepsilon_0}\ln\frac r{r_0}.
\end{equation}
よって $\phi(r_0)=0$ とすれば、原問の式を得る。

\section*{[2] 平行な二本の線電荷}
点 $(r\cos\theta,r\sin\theta)$ から $(a,0),(b,0)$ までの距離をそれぞれ $r_a,r_b$ とすると
\begin{equation}
 r_a^2=a^2+r^2-2ar\cos\theta,\qquad
 r_b^2=b^2+r^2-2br\cos\theta.
\end{equation}
[1]のポテンシャルを重ね合わせて
\begin{equation}
 \boxed{\phi(r,\theta)=\frac{\lambda}{2\pi\varepsilon_0}\ln\frac{r_b}{r_a}
 =\frac{\lambda}{4\pi\varepsilon_0}
 \ln\frac{b^2+r^2-2br\cos\theta}{a^2+r^2-2ar\cos\theta}.}
\end{equation}

\section*{[3] 鏡像電荷と境界条件}
線電荷 $+\lambda$ を $(d,0)$、鏡像線電荷 $-\lambda$ を $(D,0)$ に置くと、全ポテンシャルは
\begin{equation}
 \phi(r,\theta)=\frac{\lambda}{4\pi\varepsilon_0}
 \ln\frac{D^2+r^2-2Dr\cos\theta}{d^2+r^2-2dr\cos\theta}.
\end{equation}
$r=R$ でこの比が $\theta$ に依存しないためには、分子・分母の $\cos\theta$ 係数で割った定数項が等しくなければならない。従って
\begin{equation}
 \frac{D^2+R^2}{2DR}=\frac{d^2+R^2}{2dR},\qquad
 (D-d)(dD-R^2)=0.
\end{equation}
像電荷は空洞外にあるので $D>R>d$、従って
\begin{equation}
 \boxed{D=\frac{R^2}{d}}.
\end{equation}
この鏡像電荷対だけのポテンシャルは $r=R$ で
$\lambda\ln(R/d)/(2\pi\varepsilon_0)$ となる。接地条件 $\phi(R,\theta)=0$ を満たすため、この定数を引く。
$D=R^2/d$ を代入してまとめると、空洞内 ($r<R$) のポテンシャルは
\begin{equation}
 \boxed{\phi(r,\theta)=\frac{\lambda}{4\pi\varepsilon_0}
 \ln\left[
 \frac{d^2r^2/R^2+R^2-2dr\cos\theta}
 {d^2+r^2-2dr\cos\theta}\right].}
\end{equation}
$r=R$ では分子と分母が等しく、対数の中身が1となるため、接地条件を満たす。加えた定数は電場および誘導面電荷密度には影響しない。

\begin{center}
\begin{tikzpicture}[>=Stealth,scale=0.9]
 \def\R{1.35}
 \def\eta{0.32}
 \pgfmathsetmacro{\d}{\eta*\R}
 \pgfmathsetmacro{\D}{\R/\eta}
 \pgfmathsetmacro{\rx}{\R*cos(140)}
 \pgfmathsetmacro{\ry}{\R*sin(140)}
 \draw[dashed] (0,0) circle (\R);
 \draw[->] (-1.7,0)--({\D+0.65},0) node[right] {$x$};
 \draw[->] (0,-1.7)--(0,1.7) node[above] {$y$};
 \fill (\d,0) circle (2pt);
 \node[above] at (\d,0.12) {$+\lambda$};
 \fill (\D,0) circle (2pt);
 \node[above] at (\D,0.12) {$-\lambda$};
 \draw[<->] (0,-0.3)--(\d,-0.3) node[midway,below] {$d$};
 \draw[->] (0,0)--(\rx,\ry) node[midway,above left] {$R$};
 \node[anchor=north] at ({(\R+\D)/2},-0.28) {$D=R^2/d$};
 \node[above left] at ({1.4*\rx},{1.4*\ry}) {$r=R$};
\end{tikzpicture}
\end{center}

\section*{[4] 誘導面電荷密度}
境界では原問の法線規約により $\sigma(\theta)=\varepsilon_0\partial_r\phi|_{r=R}$ である。上のポテンシャルを微分し、$D=R^2/d$ を代入すると
\begin{align}
 \sigma(\theta)
 &=\frac{\lambda}{4\pi}\left[
 \frac{2d^2/R-2d\cos\theta}{d^2+R^2-2dR\cos\theta}
 -\frac{2R-2d\cos\theta}{d^2+R^2-2dR\cos\theta}\right]\\
 &=\boxed{-\frac{\lambda(R^2-d^2)}
 {2\pi R(R^2+d^2-2Rd\cos\theta)}}.
\end{align}
正の線電荷に対して誘起面電荷は全周で負であり、最も近い $\theta=0$ 側で絶対値が最大となる。

\section*{[5] 誘起電荷の全量}
$a=d/R<1$ と置くと
\begin{equation}
 \sigma(\theta)=-\frac{\lambda}{2\pi R}
 \frac{1-a^2}{1-2a\cos\theta+a^2}
 =-\frac{\lambda}{2\pi R}
 \left(1+2\sum_{n=1}^{\infty}a^n\cos n\theta\right).
\end{equation}
各余弦の一周積分は0なので、円筒内側面の単位長さ当たり電荷は
\begin{equation}
 \boxed{R\int_0^{2\pi}\sigma(\theta)\,d\theta=-\lambda}.
\end{equation}

\section*{[6] 面密度測定から位置を求める}
線電荷の角度が $\psi$ の場合、[4]の式は
\begin{equation}
 \sigma(\theta)=-\frac{\lambda}{2\pi R}
 \left[1+2\frac dR\cos(\theta-\psi)\right]
 +O\!\left((d/R)^2\right)
\end{equation}
となる。以下の対向点の和・差と位置の復元も $d/R$ の一次までであり、対向点で和と差を取れば定数項と反対称な成分を分離できる。
\begin{align}
 \sigma_1-\sigma_3&=-\frac{2\lambda d}{\pi R^2}\cos\psi,&
 \sigma_1+\sigma_3&=-\frac{\lambda}{\pi R},\\
 \sigma_2-\sigma_4&=-\frac{2\lambda d}{\pi R^2}\sin\psi,&
 \sigma_2+\sigma_4&=-\frac{\lambda}{\pi R}.
\end{align}
従って、$x_0=d\cos\psi$, $y_0=d\sin\psi$ より
\begin{equation}
 \boxed{x_0=\frac R2\frac{\sigma_1-\sigma_3}{\sigma_1+\sigma_3},\qquad
 y_0=\frac R2\frac{\sigma_2-\sigma_4}{\sigma_2+\sigma_4}.}
\end{equation}

\end{document}
